\section{总结与延伸}

这场 3 小时 15 分钟的访谈，核心不是“某个 Agent 项目火了”，而是展示了 Agentic Engineering 在真实世界的完整剖面：产品定义、系统实现、安全治理、组织选择与社会后果。

\subsection{核心要点总表}
\begin{longtable}{p{0.18\textwidth}p{0.3\textwidth}p{0.42\textwidth}}
\toprule
\textbf{主题} & \textbf{访谈结论} & \textbf{可执行启示} \\
\midrule
产品形态 & Agent 从对话工具进化为行动系统 & 以任务闭环而非问答质量定义产品成功 \\
\midrule
系统架构 & loop + memory + harness 构成可自治最小系统 & 先设计权限边界，再放大自治能力 \\
\midrule
工程方法 & agentic engineering 优于无约束 vibe coding & 显式区分讨论模式与执行模式 \\
\midrule
安全治理 & 风险是复合型、持续型而非一次性修补 & 常态化审计、最小权限、可追溯日志 \\
\midrule
市场变化 & app 价值向 API/数据可调用性迁移 & 尽早建设 agent-facing 接口能力 \\
\midrule
职业变化 & 手工编码价值下沉，系统设计价值上升 & 强化建模、架构、验证与协作能力 \\
\midrule
社会影响 & 增益与阵痛并存，且存在时间错位 & 技术推广需配套转型支持与教育重构 \\
\bottomrule
\end{longtable}

\subsection{部署前检查清单}
\begin{longtable}{p{0.34\textwidth}p{0.52\textwidth}}
\toprule
\textbf{检查项} & \textbf{达标标准} \\
\midrule
任务边界定义 & 每个自动化任务都有“可做/不可做”清单，且可在配置中查看。 \\
\midrule
权限分级 & L1-L4 权限模型启用，高风险动作默认二次确认。 \\
\midrule
模型路由策略 & 不同风险任务绑定不同模型，且路由规则可审计。 \\
\midrule
工具超时与重试 & 所有外部工具有超时、重试上限与熔断策略。 \\
\midrule
输入净化 & 消息、网页、文件入口均经过结构化预处理。 \\
\midrule
记忆治理 & memory 写入有 schema 校验与版本记录。 \\
\midrule
日志可追踪性 & 每次执行有 request-id，可回放关键步骤。 \\
\midrule
安全审计脚本 & 发布前自动运行审计，失败即阻断部署。 \\
\midrule
回滚机制 & 配置与策略支持一键回滚，且验证过可用。 \\
\midrule
人工接管路径 & 任意任务可在 1 步内切回人工处理。 \\
\midrule
告警策略 & 越权调用、异常高频失败、外发敏感数据均有告警。 \\
\midrule
成本控制 & token、工具调用、第三方 API 费用均有预算阈值。 \\
\midrule
评测体系 & 同时追踪 Task Success、Reliability、Safety、Business。 \\
\midrule
用户告知与授权 & 高风险能力有明确告知与用户可撤销授权。 \\
\midrule
周度复盘机制 & 固定复盘失败案例并更新规则库与测试集。 \\
\bottomrule
\end{longtable}

\subsection{进一步阅读}
\begin{itemize}
    \item OpenClaw 项目仓库与安全文档（含审计脚本与配置建议）。
    \item Lex Fridman Podcast \#491 完整视频与评论区高频问题讨论。
    \item 相关议题：Prompt Injection、防御性 Agent Runtime、Capability Sandboxing。
    \item 人机协作方法：从“写代码”到“编排执行系统”的工程实践。
\end{itemize}

\subsection{开放问题}
\begin{itemize}
    \item 如何为 memory 文件建立可验证 provenance，避免“身份污染”？
    \item 如何在保持可扩展性的同时给出细粒度权限预算（time/tool/cost）？
    \item 当 Agent 成为默认入口，平台应如何定义“用户授权代理访问”的新规范？
\end{itemize}

\end{document}
